\section*{Hoofdstuk \ref{ch:b3:bioinorganic} --- Bio-anorganische chemie}

\begin{solution}{exo:b3:bioinorganic:1}
$\log(\theta/(1 - \theta))$ gaat van $\log 0.25 = -0.602$ naar $\log 4 = +0.602$ terwijl $\log p$ stijgt met
$\log(5.9/2.1) = 0.449$: $n = 1.204/0.449 = 2.7$.
\end{solution}

\begin{solution}{exo:b3:bioinorganic:2}
$\ce{Mn^2+} < \ce{Fe^2+} < \ce{Co^2+} < \ce{Ni^2+} < \ce{Cu^2+} > \ce{Zn^2+}$. Zink ($d^{10}$) heeft geen ligandveldstabilisatie en geen
vervorming van Jahn-Teller: het valt terug onder koper.
\end{solution}

\begin{solution}{exo:b3:bioinorganic:3}
$\ce{Ca^2+}$ en $\ce{Fe^3+}$, hard: carboxylaten (en fenolaten voor ijzer). $\ce{Cu+}$ en $\ce{Hg^2+}$, zacht:
cysteïnethiolaten. $\ce{Zn^2+}$, grensgeval: histidine en cysteïne.
\end{solution}

\begin{solution}{exo:b3:bioinorganic:4}
$\ce{2H2O -> O2 + 4H+ + 4e-}$: vier elektronen, vier fotonen per $\ce{O2}$; vier mol fotonen per
mol $\ce{O2}$.
\end{solution}

\begin{solution}{exo:b3:bioinorganic:5}
Myoglobine: $\theta = 13/13.37 = 0.972$ en $5/5.37 = 0.931$: het staat $(0.972 - 0.931)/0.972 = 4$\,\% van zijn
zuurstof af. Hemoglobine: $\theta = 0.972$ en $0.724$: het staat 26\,\% af.
\end{solution}

\begin{solution}{exo:b3:bioinorganic:6}
$[\mathrm{HbCO}]/[\mathrm{HbO_2}] = 230 \times \num{50e-6}/0.2095 = 0.055$; fractie $0.055/1.055 = 5.2$\,\%.
\end{solution}

\begin{solution}{exo:b3:bioinorganic:7}
Sulfiden elk $-2$, cysteïnaten elk $-1$. $\ce{[Fe2S2(SR)4]^2-}$: de ijzerionen samen $-2 + 4 + 4 = +6$: twee
Fe(III). $\ce{[Fe2S2(SR)4]^3-}$: $+5$, één Fe(III) en één Fe(II) (gemengde valentie). $\ce{[Fe4S4(SR)4]^2-}$:
$-2 + 8 + 4 = +10$, twee Fe(III) en twee Fe(II), gemiddeld $+2.5$ per ijzerion.
\end{solution}

\begin{solution}{exo:b3:bioinorganic:8}
$\ce{[PtCl2(NH3)2] + H2O -> [PtCl(H2O)(NH3)2]+ + Cl-}$. De hoge chlorideconcentratie van het bloed
duwt het evenwicht terug; in cellen is de chlorideconcentratie veel
lager, en ontstaat het aquacomplex. Platina bindt aan N7 van guanine. In
het isomeer \textit{trans} liggen de twee
labiele plaatsen tegenover elkaar en kunnen ze geen twee naburige basen
op één streng bereiken.
\end{solution}

\begin{solution}{exo:b3:bioinorganic:9}
$1/T_1 = 1/\qty{1.2}{s} + \qty{4.0}{L.mmol^{-1}.s^{-1}} \times \qty{0.10}{mmol/L} = 0.833 + 0.400 = \qty{1.23}{s^{-1}}$: $T_1 = \qty{0.81}{s}$.
\end{solution}

\begin{solution}{exo:b3:bioinorganic:10}
$K = K_{\mathrm{PbL}}/K_{\mathrm{CaL}} = 10^{18.0 - 10.7} = 10^{7.3} = \num{2e7}$: lood verdringt calcium volledig. Het vrije
ligand zou ook het calcium van het bloed binden en de concentratie ervan
gevaarlijk verlagen; het calciumcomplex wisselt uit met lood en laat het
calciumniveau ongewijzigd.
\end{solution}

\begin{solution}{exo:b3:bioinorganic:11}
Ver onder $K_M$ is de reactie van de eerste orde: $k' = (k_{\mathrm{cat}}/K_M)[\mathrm E]$, $t_{1/2} = \ln 2/k'$. Koolzuuranhydrase:
$k' = \qty{100}{s^{-1}}$, $t_{1/2} = \qty{6.9}{ms}$. Typisch \omterm{def:b3:complex-kinetics:enzyme}{enzym}: $k' = \qty{0.10}{s^{-1}}$, $t_{1/2} = \qty{6.9}{s}$, duizend
keer langer, te traag voor de seconde die een rode bloedcel in een
longcapillair doorbrengt.
\end{solution}

\begin{solution}{exo:b3:bioinorganic:12}
De helft van de plaatsen met CO betekent $[\mathrm{HbCO}] = [\mathrm{HbO_2}]$: $p(\ce{CO}) = p(\ce{O2})/M = 0.2095/230 = \qty{9.1e-4}{bar}$, ongeveer
\qty{910}{ppm}.
\end{solution}

\begin{solution}{pb:b3:bioinorganic:1}
\textbf{1.} $\theta = 13^{2.7}/(3.5^{2.7} + 13^{2.7}) = 0.972$.

\textbf{2.} $\theta(\qty{5.0}{kPa}) = 0.724$.

\textbf{3.} $5.0/(0.37 + 5.0) = 0.931$.

\textbf{4.} Zijn $p_{50}$ is tien keer lager: bij de drukken van een werkende spier is het veel sterker
verzadigd dan hemoglobine bij dezelfde druk, zodat zuurstof van
hemoglobine naar myoglobine overgaat.

\textbf{5.} $n = 1$.

\textbf{6.} Binding aan één \omterm{def:b3:bioinorganic:haem}{heem} verschuift zijn ijzer in het vlak en
trekt de histidine en haar helix mee; de subeenheden verschuiven samen
naar een vorm met een hogere affiniteit voor de andere plaatsen.

\textbf{7.} $\qty{150}{g/L}/\qty{64500}{g/mol} = \qty{2.33e-3}{mol/L}$.

\textbf{8.} $4 \times \num{2.33e-3} = \qty{9.30e-3}{mol/L}$.

\textbf{9.} $\qty{9.30}{mmol/L} \times 0.972 = \qty{9.04}{mmol}$.

\textbf{10.} $\qty{9.30}{mmol/L} \times (0.972 - 0.724) = \qty{2.31}{mmol}$.

\textbf{11.} $\qty{2.31e-3}{mol} \times \qty{22.41}{L/mol} = \qty{52}{mL}$.

\textbf{12.} $0.248/0.972 = 26$\,\%.

\textbf{13.} $5.0^{2.7}/(4.0^{2.7} + 5.0^{2.7}) = 0.646$.

\textbf{14.} $\qty{9.30}{mmol/L} \times (0.972 - 0.646) = \qty{3.03}{mmol}$.

\textbf{15.} $3.03/2.31 = 1.31$: dertig procent meer.

\textbf{16.} Koolstofdioxide uit de ademhaling, door koolzuuranhydrase
gehydrateerd tot koolzuur, en melkzuur bij zware inspanning.

\textbf{17.} $\qty{3.03}{mmol/L} \times \qty{5.0}{L/min} = \qty{15.1}{mmol/min}$.

\textbf{18.} $\qty{15.1e-3}{mol/min} \times \qty{22.41}{L/mol} = \qty{0.34}{L/min}$.

\textbf{19.} $230 \times \num{100e-6}/0.2095 = 0.110$.

\textbf{20.} $0.110/1.110 = 0.099$: ongeveer 10\,\% van de plaatsen.

\textbf{21.} De plaatsen die nog zuurstof dragen, in een molecuul dat ook
CO draagt, zijn in de vorm met hoge affiniteit: de kromme verschuift naar
links en ze geven in de weefsels minder van hun zuurstof af.

\textbf{22.} Volgens de competitierelatie daalt de verhouding $[\mathrm{HbCO}]/[\mathrm{HbO_2}]$ evenredig met $p(\ce{O2})$:
zuivere zuurstof inademen, bijna vijf keer de druk in lucht, verdringt
het koolstofmonoxide sneller.

\textbf{23.} Door de lineaire binding van CO te hinderen en een
waterstofbrug te vormen met $\ce{O2}$, houdt ze $M$ in
de honderden in plaats van de veel grotere waarde van vrij \omterm{def:b3:bioinorganic:haem}{heem}; anders
zou zelfs het koolstofmonoxide dat het lichaam zelf maakt, een groot deel
van het hemoglobine blokkeren.

\textbf{24.} In rust levert het bloed ongeveer \qty{0.34}{L} zuurstof per minuut.
\end{solution}
